\section*{अध्याय \ref{ch:g10:lines} --- रेखाओं के समीकरण और रैखिक निकाय}

\begin{solution}{exo:g10:lines:1}
$y = 2x - 3$: \omterm{def:g10:reffunc:affine}{प्रवणता} $2$, अंतःखंड $-3$।
$y = -\frac13 x + 1$: \omterm{def:g10:reffunc:affine}{प्रवणता} $-\frac13$, अंतःखंड $1$।
$y = 4$: \omterm{def:g10:reffunc:affine}{प्रवणता} $0$, अंतःखंड $4$ (क्षैतिज रेखा)।
$x = -2$: ऊर्ध्वाधर रेखा, न \omterm{def:g10:reffunc:affine}{प्रवणता}, न \omterm{thm:g10:lines:reduced}{संक्षिप्त समीकरण}।
\end{solution}

\begin{solution}{exo:g10:lines:2}
$3 \times 2 + 1 = 7$: हाँ, $A$ रेखा पर है।
$3 \times (-1) + 1 = -2 \neq -1$: $B$ नहीं है।
$3 \times 4 + 1 = 13$: हाँ, $C$ रेखा पर है।
\end{solution}

\begin{solution}{exo:g10:lines:3}
(a) $m = \frac{8 - 2}{3 - 0} = 2$ और $p = y_A = 2$ (बिंदु $A$ $y$-अक्ष पर है): $y = 2x + 2$।

(b) $m = \frac{-4 - 5}{2 - (-1)} = \frac{-9}{3} = -3$; फिर $5 = -3 \times (-1) + p$ से $p = 2$ मिलता है: $y = -3x + 2$।

(c) $x_A = x_B = 4$: ऊर्ध्वाधर रेखा $x = 4$।
\end{solution}

\begin{solution}{exo:g10:lines:4}
$y = \frac{6x+2}{2} = 3x + 1$। \omterm{def:g10:reffunc:affine}{प्रवणता} $3$ वाली रेखाएँ $y = 3x - 1$, $y = 3x + 5$ और $y = 3x + 1$ हैं: ये तीनों परस्पर
समांतर हैं; इनमें से कोई दो बराबर नहीं हैं (अंतःखंड भिन्न हैं), और इनमें से
कोई $y = -3x + 2$ के समांतर नहीं है।
\end{solution}

\begin{solution}{exo:g10:lines:5}
पहला निकाय: $3x + y = 9$ में $y = 2x - 1$ रखिए: $3x + 2x - 1 = 9$, इसलिए $5x = 10$, $x = 2$, फिर $y = 3$। हल: $(2, 3)$।

दूसरा निकाय: $x - y = 4$ से $x = y + 4$। तब $2(y + 4) + 3y = 3$, इसलिए $5y = -5$, $y = -1$, फिर $x = 3$। हल: $(3, -1)$।
\end{solution}

\begin{solution}{exo:g10:lines:6}
पहला निकाय: \omterm{def:g10:algebra:equation}{समीकरणों} को जोड़ने पर $y$ का विलोप हो जाता है: $2x = 14$, $x = 7$; फिर
$7 + y = 10$ से $y = 3$ मिलता है। हल: $(7, 3)$।

दूसरा निकाय: पहले \omterm{def:g10:algebra:equation}{समीकरण} को $5$ से और दूसरे को $-2$ से गुणा कीजिए:
\[
\begin{cases}
15x + 10y = 60 \\
-4x - 10y = -16
\end{cases}
\]
जोड़ने पर: $11x = 44$, इसलिए $x = 4$; फिर $3 \times 4 + 2y = 12$ से $y = 0$ मिलता है। हल: $(4, 0)$।
\end{solution}

\begin{solution}{exo:g10:lines:7}
हर स्थिति में $ab' - a'b$ निकालिए।

पहली: $2 \times 2 - (-4) \times (-1) = 4 - 4 = 0$, और दूसरा \omterm{def:g10:algebra:equation}{समीकरण} पहले का $-2$ गुना है: एक ही रेखा दो बार, अनंत
हल।

दूसरी: फिर $ab' - a'b = 0$, पर दाहिने पक्ष $-2$ के अनुपात में नहीं हैं ($1 \neq -6$): दो भिन्न
समांतर रेखाएँ, कोई हल नहीं।

तीसरी: $2 \times 1 - 1 \times (-1) = 3 \neq 0$: ठीक एक हल।
\end{solution}

\begin{solution}{exo:g10:lines:8}
मान लीजिए $x$ बच्चों के टिकटों की संख्या है और $y$ वयस्कों के टिकटों की:
\[
\begin{cases}
x + y = 200 \\
5x + 9y = 1352 .
\end{cases}
\]
पहले \omterm{def:g10:algebra:equation}{समीकरण} से $x = 200 - y$; प्रतिस्थापन करने पर $5(200 - y) + 9y = 1352$, इसलिए $1000 + 4y = 1352$, $y = 88$, फिर $x = 112$।
अतः $112$ बच्चों के और $88$ वयस्कों के टिकट। जाँच: $5 \times 112 + 9 \times 88 = 560 + 792 = 1352$।
\end{solution}

\begin{solution}{exo:g10:lines:9}
\emph{1.} समांतर होने का अर्थ है वही \omterm{def:g10:reffunc:affine}{प्रवणता} $2$: $y = 2x + p$, जहाँ $4 = 2 \times 1 + p$, इसलिए
$p = 2$: $d'\colon y = 2x + 2$।

\emph{2.} $x$-अक्ष पर $y = 0$: $2x + 2 = 0$ से $x = -1$ मिलता है। प्रतिच्छेद बिंदु $(-1, 0)$
है।
\end{solution}

\begin{solution}{exo:g10:lines:10}
\emph{1.} $[BC]$ का \omterm{prop:g10:coordgeom:midpoint}{मध्यबिंदु}: $A' = \left(\frac{6+2}{2},
\frac{2+4}{2}\right) = (4, 3)$। $A(0,0)$ से खींची माध्यिका की \omterm{def:g10:reffunc:affine}{प्रवणता}
$\frac{3 - 0}{4 - 0} = \frac34$ है: \omterm{def:g10:algebra:equation}{समीकरण} $y = \frac34 x$।

$[AC]$ का \omterm{prop:g10:coordgeom:midpoint}{मध्यबिंदु}: $B' = (1, 2)$। $B(6, 2)$ से खींची माध्यिका की \omterm{def:g10:reffunc:affine}{प्रवणता} $\frac{2 - 2}{1 - 6} = 0$ है: वह
क्षैतिज रेखा $y = 2$ है।

\emph{2.} प्रतिच्छेद: $\frac34 x = 2$ से $x = \frac83$ मिलता है, इसलिए $G\left(\frac83, 2\right)$। तीसरी माध्यिका
$C(2,4)$ को $[AB]$ के \omterm{prop:g10:coordgeom:midpoint}{मध्यबिंदु} $C' = (3, 1)$ से जोड़ती है; उसकी \omterm{def:g10:reffunc:affine}{प्रवणता} $\frac{1 - 4}{3 - 2} = -3$ है,
\omterm{def:g10:algebra:equation}{समीकरण} $y = -3x + 10$। $x = \frac83$ पर: $y = -8 + 10 = 2$: $G$ उस पर है। और सचमुच $G = \left(\frac{0 + 6 + 2}{3}, \frac{0 + 2 + 4}{3}\right)$: केन्द्रक
शीर्षों के \omterm{def:g10:coordgeom:system}{निर्देशांकों} का माध्य ले लेता है।
\end{solution}

\begin{solution}{pb:g10:lines:1}
\textbf{1.} \omterm{def:g10:reffunc:affine}{प्रवणता} $\frac{8 - 2}{3 - 1} = 3$; $(1, 2)$ से होकर: $y = 3x - 1$।

\textbf{2.} $y = 3x - 1$ और $y = 3x + 4$: वही \omterm{def:g10:reffunc:affine}{प्रवणता}, समांतर (और भिन्न)। $y = 3x - 1$ को $y = -x + 7$
से काटने पर: $3x - 1 = -x + 7$ से $x = 2$, $y = 5$ मिलता है: बिंदु $(2, 5)$।

\textbf{3.} दोनों बिंदुओं का $x = 2$ एक ही है: रेखा ऊर्ध्वाधर है, \omterm{def:g10:algebra:equation}{समीकरण}
$x = 2$। संक्षिप्त रूप $y = mx + p$ इस प्रश्न का उत्तर देता है कि ``प्रति $x$
कितना $y$'', और ऊर्ध्वाधर रेखा यह मानने से मना कर देती है: उसकी \omterm{def:g10:reffunc:affine}{प्रवणता}
अनंत होगी।

\textbf{4.} $3 \times 4 - 1 = 11$: हाँ, $(4, 11)$ रेखा पर है। $3 \times (-1) - 1 = -4 \neq -5$: नहीं।

\textbf{5.} $y = -2x + 3$; अक्षों को $(0, 3)$ और $\left(\frac32, 0\right)$ पर काटती है। त्रिभुज का
क्षेत्रफल: $\frac12 \times \frac32 \times 3 = \frac94 = 2.25$।

\textbf{6.} प्रतिस्थापन करने पर: $3x + 2x - 3 = 7$, $x = 2$, $y = 1$: हल $(2, 1)$।

\textbf{7.} जोड़ने पर: $7x = 7$, $x = 1$; फिर $3y = 6$, $y = 2$: हल $(1, 2)$।

\textbf{8.} पहला निकाय: दूसरा \omterm{def:g10:algebra:equation}{समीकरण} पहले का दुगुना है: एक ही रेखा दो बार
गिनी गई --- अनंत हल ($2x - y = 3$ के सभी बिंदु)। दूसरा निकाय: बाएँ पक्ष उसी अनुपात
में हैं, दाहिने नहीं ($2 \times 3 = 6 \neq 1$): दो समांतर रेखाएँ --- कोई हल नहीं। त्रिविभाजन:
भिन्न प्रवणताएँ $\to$ एक प्रतिच्छेद; समान \omterm{def:g10:reffunc:affine}{प्रवणता}, भिन्न अंतःखंड $\to$
समांतर, कोई हल नहीं; सब कुछ समान $\to$ एक ही रेखा, अनंत हल।

\textbf{9.} रेखाओं के दिक् \omterm{def:g10:vectors:vector}{सदिश} $(-b, a)$ और $(-d, c)$ हैं (या: प्रवणताएँ $-\frac ab$ और
$-\frac cd$), और वे \omterm{def:g10:vectors:collinear}{संरेखी} ठीक तब हैं जब $ad - bc = 0$ हो --- \cref{pb:g10:vectors:1} वाला सारणिक नई नौकरी
पर। मान: प्रश्न 7: $2 \times (-3) - 3 \times 5 = -21 \neq 0$: अद्वितीय हल। प्रश्न 8: दोनों के लिए $2 \times (-2) - (-1) \times 4 = 0$:
समांतर-या-समरूप, और अचर पद दोनों स्थितियों को अलग कर देते हैं।

\textbf{10.} $ad - bc = 6 - 2m$: $m = 3$ के लिए शून्य। वहाँ $x + 3y = 3$ के मुक़ाबले $2x + 6y = 7$: पहले
को दुगुना करने पर $2x + 6y = 6 \neq 7$ मिलता है: समांतर, \emph{कोई हल नहीं}। $m \neq 3$ के लिए
विलोपन अद्वितीय हल $y = \frac{1}{6 - 2m}$, फिर $x = 3 - \frac{m}{6 - 2m}$ देता है।

\textbf{11.} $p + r = 35$, $2p + 4r = 94$: आधा करने पर $p + 2r = 47$, इसलिए $r = 12$ और $p = 23$: तेईस
तीतर, बारह ख़रगोश --- \emph{नौ अध्याय} का अपना उत्तर।

\textbf{12.} $x + y = 30$ और $0.6x + 0.2y = 0.35 \times 30 =
10.5$: पहले में से $0.2 \times$ घटाने पर $0.4x = 4.5$: $x = 11.25$ L शरबत,
$y = 18.75$ L पेय।

\textbf{13.} अंक $a$, $b$: $a + b = 12$ और $(10a + b) - (10b + a) = 9(a - b) = 18$, इसलिए $a - b = 2$: $a = 7$, $b = 5$:
संख्या $75$ है।

\textbf{14.} जोड़ने पर: $5c + 5r = 18$, इसलिए $c + r = 3.60$ --- एक कॉफ़ी और एक क्रोसां।
घटाने पर: $c - r = 0.20$। अतः $c = 1.90$ और $r = 1.70$ यूरो: सममित छोटे रास्ते ने बिना किसी
प्रतिस्थापन के हल कर दिया।

\textbf{15.} पहले दावे को दुगुना करने पर $4$ सेब $+ 2$ नाशपाती $= 10$ यूरो
में मिलनी चाहिए, पर दुकान $9$ लेती है: निकाय असंगत है --- समांतर रेखाएँ,
रिक्त हल-समुच्चय। निर्णय: दोनों दाम एक साथ सही नहीं हो सकते; या तो कोई छूट
छिपी है, या हिसाब।

\textbf{16.} मूल बिंदु पर: $0 > 3 \times 0 - 1 = -1$: सत्य, इसलिए $(0,0)$ क्षेत्र $y > 3x - 1$ में है
(रेखा के ऊपर)। किसी भी बिंदु को रखकर जाँच लीजिए: असमिका सत्य निकले तो ऊपर,
नहीं तो नीचे।

\textbf{17.} कोने: $(0, 0)$; $(4, 0)$ (अक्ष और $x + y = 4$); $(0, 1)$ (अक्ष और $y = 2x + 1$); और
$x + y = 4$ तथा $y = 2x + 1$: $x = 1$, $y = 3$: $(1, 3)$।

\textbf{18.} $P(0,0) = 0$; $P(4, 0) = 12$; $P(0, 1) = 2$; $P(1, 3) = 9$। \omterm{def:g10:functions:extrema}{अधिकतम} कोने $(4, 0)$ पर: योजना $x = 4$,
$y = 0$ से $12$ मिलता है --- कोनों का राज, जैसा बहती समांतर रेखाएँ $3x + 2y = c$
दिखा देती हैं।

\textbf{19.} $1.5x + 3 = 2x$ से $x = 6$, $y = 12$ मिलता है: छह किलोमीटर पर दोनों टैक्सियाँ
बारह यूरो लेती हैं --- माध्यमिक विद्यालय खंड की दो-थर्मामीटर वाली सप्ताहांत
समस्या का वही बराबरी वाला बिंदु, अब किसी निकाय का अद्वितीय हल।

\textbf{20.} किसी रैखिक \omterm{def:g10:algebra:equation}{समीकरण} के हल एक रेखा खींचते हैं; निकाय दो रेखाएँ एक
साथ रख देता है, और $ad - bc$ एक नज़र में बता देता है कि वे एक बार कटती हैं
(शून्येतर) या समांतर-या-समरूप वाली स्थितियों में गिरती हैं (शून्य)। इसके
बदले असमिकाएँ ढेर कर दीजिए तो हल एक बहुभुजीय क्षेत्र बनाते हैं जिसके कोने
हर रैखिक इष्टतम उठाए रहते हैं। रैखिक बीजगणित सारणिक को किसी भी संख्या के
अज्ञातों तक सामान्य कर देता है; रैखिक प्रोग्रामन कोनों को उद्योग बना देता
है।
\end{solution}
